\section{结论}

本文提出一套把公开多工具任务转换为冗余参数泄露评测数据的可复现方法。该方法将来源解析、来源策略、统一筛选、隐私构造、独立验证和来源重放分开，每个步骤的输出都可以由校验器从发布文件重新计算。泄露调用中每个新增参数都与一个档案字段同名、同值，功能调用则完整保留上游参考参数，两者的差异精确落在新增的可选敏感参数上。

基于锁定的 $\tau$-bench、BFCL、API-Bank 和 AppWorld，本文构建 328 条任务，覆盖客服、交易、工具检索和跨应用执行四种任务形态。全部任务通过结构校验和来源重放，其中 321 条通过最终状态验证。从四个来源分层抽取 80 条授权关系进行人工复核，78 条与确定性标签一致，2 条不一致来自值匹配规则在语义边界上的局限。

本文的主要贡献在于提出了一种主动构造配对样本的方法：不依赖智能体的自然执行行为来观察泄露，而是为每个有效步骤预先构造功能相同、披露程度不同的两个调用版本。这使得泄露差异具有确定的结构，可以反复校验和复现，也为后续的参数级隐私评测提供了可执行的数据基础。

后续改进包括：分析原生可选参数而非合成扩展的泄露机会；开展参数消融实验以评估参考调用的参数最小性；完成两名独立复核者的人工双审并计算一致性指标。
